% answers_short.tex —— 2009 年·普通高等学校招生考试·四川卷·理科 · 简短答案速查
% 由 tools/make_short_answers.py 自 answers.tex 的 \gksolhead/\gkanswer 提取，
% 并与 2009/2009_数学真题_四川理_答案要点.md 逐题比对后生成，勿手改（改源 MD 后重跑）。

\begin{gknotice}
\textbf{使用说明}：本节只给最终答案便于快速核对，完整推导见解析版。答案为编写者独立推导，\textbf{非官方标准答案}。
\end{gknotice}

\begin{gknotice}
选择题答案中的 U 表示待核，未给出猜测答案，详情见解析版。
\end{gknotice}

\gkpart{选择题}
\begin{center}\begin{tabular}{ccccccccccccc}
\toprule
题号 & 1 & 2 & 3 & 4 & 5 & 6 & 7 & 8 & 9 & 10 & 11 & 12 \\
\midrule
答案 & C & B & A & D & D & B & C & B & A & D & U & A \\
\bottomrule
\end{tabular}\end{center}

\gkpart{填空题}
\begin{itemize}[leftmargin=2.6em]
  \item[\textbf{13.}] $-20$。
  \item[\textbf{14.}] $4$。
  \item[\textbf{15.}] $\frac{\pi}{2}$。
  \item[\textbf{16.}] ①②④。
\end{itemize}

\gkpart{解答题}
\begin{itemize}[leftmargin=2.6em]
  \item[\textbf{17.}] （1）$A+B=\frac\pi4$；（2）$a=\sqrt2$，$b=1$，$c=\sqrt5$。
  \item[\textbf{18.}] （1）$\frac{36}{85}$；（2）分布列见下，$E\xi=2$。
  \item[\textbf{19.}] （1）证明见解析；（2）存在，$M$ 为 $AE$ 的中点；（3）$\arccos\frac{3\sqrt{11}}{11}$。
  \item[\textbf{20.}] （1）$\frac{x^2}{2}+y^2=1$；（2）$y=x+1$ 或 $y=-x-1$。
  \item[\textbf{21.}] （1）$0<a<1$ 时定义域 $(0,+\infty)$，$a>1$ 时 $(-\infty,0)$，两种情形下均单调递减；（2）$\frac1a$；（3）$m>0$，极大值 $2(1+\sqrt m)\mathrm e^{-1-\sqrt m}$，且 $0<m<1$ 时极小值为 $2(1-\sqrt m)\mathrm e^{-1+\sqrt m}$。
  \item[\textbf{22.}] （1）$b_n=4+\dfrac{5(-1)^n}{4^n-(-1)^n}$；（2）证明见解析；（3）$\lambda$ 的最小值为 $4$。
\end{itemize}
